The sixth example demonstrates the image-generation pathway, in which the
pipeline intercepts a creative image prompt, identifies a matching sponsored
product, and rewrites the prompt so that the product appears naturally within
the generated scene.  Figure~\ref{fig:example-6} shows the result: a
photorealistic sunrise trail image in which the runner is visibly wearing the
Nike Pegasus~41.

\begin{figure}[H]
    \centering
    \includegraphics[width=0.6\textwidth]{images/illustrative_example_6.png}
    \caption{Image generated from the rewritten prompt, placing the Nike
    Pegasus~41 on the runner's feet at sunrise on a mountain trail.
    The \textsc{[Sponsored]} label and product link are returned alongside
    the image in the API response.}
    \label{fig:example-6}
\end{figure}

\subsubsection{Stages 1--2: Context Extraction and Ad Matching}
\label{subsubsec:example6-context}

The user submits a short creative prompt:
\emph{``A runner training at sunrise on a mountain trail''}.
Context extraction and vector search proceed as described in
Section~\ref{subsubsec:example1-context} and~\ref{subsubsec:example1-retrieval}.
The key difference is that the \texttt{intent} field is set to
\texttt{image\_generation}, routing the request to the image-generation handler
rather than the text-synthesis handler.  The Nike Pegasus~41 again ranks first
(similarity~$= 0.8124$) and passes the budget gate.  Notably, its inventory
record includes a \texttt{product\_image\_url} field, which signals to the
handler that an image-edit call is available as a grounded alternative to a
plain text-to-image call.

\subsubsection{Stage 3: Prompt Rewriting}
\label{subsubsec:example6-prompt}

The image-generation handler invokes the prompt-rewriter module, which
receives the user's original prompt and the winning ad's structured fields.
The rewriter's task is to produce an enhanced prompt that retains the user's
intended scene while incorporating the sponsored product in a visually
natural, diegetic manner.  The rewritten prompt returned for this request is:

\begin{quote}
\emph{A runner laces up their Nike Pegasus~41 shoes, preparing for a training
session at sunrise on a mountain trail.  The shoes are a perfect fit for this
daily run, with the lightweight design and responsive ZoomX foam ready to
tackle the easy to moderate miles ahead.  As the runner takes their first
steps, the warm sunrise light casts a glow on the trail, illuminating the
winding path and the majestic mountains in the background.  The runner's
breath is visible in the cool morning air as they begin their journey, the
Nike Pegasus~41 shoes making every step feel effortless and natural.  The
trail stretches out before them, inviting the runner to explore and push their
limits.}
\end{quote}

\noindent
The rewriter inserts the product name twice, connects the product's key
features (ZoomX foam, lightweight design) to the scene's physical context,
and frames both references as incidental details rather than promotional
claims---consistent with the diegetic-placement objective.

\subsubsection{Stage 4: Image Generation and Impression Logging}
\label{subsubsec:example6-response}

The rewritten prompt is forwarded to the image-generation endpoint.  When a
\texttt{product\_image\_url} is available the handler issues an
\texttt{images.edit} call, supplying the product photo as a reference anchor
so that the model can ground the shoe's appearance on the actual product
rather than a hallucinated approximation.  The endpoint returns a base64-encoded
image delivered as a data URI, accompanied by the ad metadata (product name,
tracking URL, and disclosure label) so that the frontend can render the
\textsc{[Sponsored]} overlay alongside the image.

Once the image is returned, the orchestrator fires the same two asynchronous
writes used in the text pathway: an impression event appended to the
\texttt{events} table, and a \texttt{daily\_spend\_log} upsert that
increments \texttt{spend} by $\text{bid\_cpm}/1000 = \$0.005$.  If no ad
had passed the similarity threshold, the handler would have fallen back to a
plain \texttt{images.generate} call with the user's original, unmodified
prompt, producing an unsponsored image with no impression recorded.
